\clearpage\begingroup\NoSourceHead
\makeatletter\tagsleft@true\let\veqno\@@leqno\makeatother
\renewcommand{\thefootnote}{\arabic{footnote})}\setcounter{footnote}{5}
\ArithmeticHead{7}{R}{GESCHICHTE DER EXTREMA VON DOPPELINTEGRALEN VOR GAUSS.}
\noindent\hspace*{\parindent}3. Obwohl \textsc{Gauss} sich der Methode der Variation der unabhängigen
Variabeln bedient, so ist es der Vollständigkeit halber doch erforderlich,
einiges über die Methode der Variation der Grenzen vorauszuschicken.

Der Grundgedanke der Methode kommt schon bei \textsc{Lagrange}\footnote{Oeuvres de \textsc{Lagrange}, Bd. XIV, S. 147—149.} vor in
einem Brief an \textsc{Euler} vom 20. November 1755 und zwar bei der Lösung des
Problems der Brachistochrone von einem gegebenen Punkt nach einer gegebenen
Kurve.

Aber erst siebzig Jahre später (1825) ist die Methode von \textsc{Ohm}\footnote{\textsc{Martin Ohm}, \textit{Die Lehre vom Grössten und Kleinsten}, Berlin 1825.} systematisch
durchgeführt und auch auf Doppelintegrale angewandt worden. Dabei
ist zunächst die Sorgfalt hervorzuheben, die \textsc{Ohm} im Gegensatz zu seinen Vorgängern
auf die Bezeichnung verwendet. So hat er bereits — siebzehn Jahre
vor \textsc{Sarrus} — ein einfaches und ein doppeltes \textls[100]{Substitutionszeichen}, und
die Grenzen der bestimmten Integrale werden stets angegeben, ebenso die
\par\clearpage\SourceLeaf{465}
\ArithmeticHead{8}{L}{BOLZA, GAUSS UND DIE VARIATIONSRECHNUNG.}
\noindent Integrationsordnung bei Doppelintegralen. Dieser sorgfältigen Bezeichnung
ist es nun auch zum grossen Teil zuzuschreiben, dass \textsc{Ohm} bei der Variation
der Doppelintegrale ganz wesentlich über seine Vorgänger \textsc{Euler} und \textsc{Lagrange}
hinausgekommen ist.

\textsc{Ohm} legt seinen Untersuchungen die \textsc{Euler}sche Auffassung der Variation
als Differentiation nach einem Parameter zugrunde. Dem entsprechend werden
in dem zu variierenden Doppelintegral
\begin{equation*}
J=\int_{x_0}^{x_1}\left(\int_{y_0(x)}^{y_1(x)}V\left(x,y,z,\frac{\partial z}{\partial x},\frac{\partial z}{\partial y},\ldots\right)dy\right)dx\tag{8}
\end{equation*}
die Grössen:
\[
z(x,y),\quad y_0(x),\quad y_1(x),\quad x_0,\quad x_1
\]
durch Funktionen eines Parameters $\varkappa$:
\[
\bar z(x,y,\varkappa),\quad\bar y_0(x,\varkappa),\quad\bar y_1(x,\varkappa),\quad\bar x_0(\varkappa),\quad\bar x_1(\varkappa)
\]
ersetzt, die sich für $\varkappa=0$ auf die vorigen reduzieren. Das so entstehende
Integral $\bar J$ wird dann nach $\varkappa$ differentiiert und dann $\varkappa=0$ gesetzt. Durch
zweimalige Anwendung der Regel für die Differentiation eines bestimmten
Integrals nach einem Parameter ergibt sich die folgende \textls[100]{erste Fundamentalformel
von \textsc{Ohm}}\footnote{\textsc{Ohm}, loc. cit. 7), S. 121, 122. \textsc{Ohm} selbst leitet die Formel (9) durch zweimalige Anwendung der von ihm vorher bewiesenen entsprechenden Formel für einfache Integrale ab.}:
\begin{equation*}
\begin{aligned}
\delta J={}&\int_{x_0}^{x_1}\int_{y_0}^{y_1}\delta V\,dy\,dx
+\int_{x_0}^{x_1}\delta y_1\Big|^{y_1}V\,dx-\int_{x_0}^{x_1}\delta y_0\Big|^{y_0}V\,dx\\
&+\delta x_1\Big|^{x_1}\int_{y_0}^{y_1}V\,dy-\delta x_0\Big|^{x_0}\int_{y_0}^{y_1}V\,dy.
\end{aligned}\tag{9}
\end{equation*}

Dabei haben wir von dem einfachen Substitutionssymbol in der aus dem
\textsc{Lindelöf-Moigno}schen Lehrbuch\begingroup\renewcommand{\thefootnote}{8a)}\footnote{\textsc{F. Moigno, L. Lindelöf}, \textit{Leçons sur le calcul différentiel et intégral}, IV, Calcul des Variations, Paris 1861.}\endgroup\setcounter{footnote}{8}\ 
bekannten Form Gebrauch gemacht, und
das Symbol $\delta$ bedeutet allgemein die Operation
\[
\delta\varphi=\left(\frac{\partial\bar\varphi}{\partial\varkappa}\right)_{\varkappa=0}.
\]
\par\clearpage\SourceLeaf{466}
\ArithmeticHead{9}{R}{GESCHICHTE DER EXTREMA VON DOPPELINTEGRALEN VOR GAUSS.}
\noindent Mit der Formel (9) war die Aufgabe der \textls[100]{Aufstellung} der ersten Variation
für Doppelintegrale mit variablen Grenzen gelöst.

4. Von fast noch grösserer Wichtigkeit ist die \textls[100]{Umformung}, die nunmehr
\textsc{Ohm} mit dem in der Formel (9) auftretenden Doppelintegral vornimmt,
nachdem er dasselbe zunächst auf die Form (3) gebracht hat. — Die Schwierigkeit,
aus der Gleichung (7) weitere Folgerungen zu ziehen, hatte im wesentlichen
darin bestanden, dass in den beiden einfachen Integralen auf der linken
Seite nach verschiedenen Variabeln und zwischen verschiedenen Grenzen integriert
wird. Um diesem Übelstand abzuhelfen, macht \textsc{Ohm} von der folgenden
Formel Gebrauch, die sich unmittelbar aus der Regel für die Differentiation
eines bestimmten Integrals nach einem Parameter ergibt,
\begin{equation*}
\int_{y_0}^{y_1}\frac{\partial A(x,y)}{\partial x}\,dy
=\frac{d}{dx}\int_{y_0}^{y_1}A(x,y)\,dy-A(x,y_1)\frac{dy_1}{dx}+A(x,y_0)\frac{dy_0}{dx}.\tag{10}
\end{equation*}
Hiernach erhält man an Stelle der \textsc{Euler}schen Formel (5) die folgende
\begin{equation*}
\int_{x_0}^{x_1}\int_{y_0}^{y_1}\left(\frac{\partial A}{\partial x}+\frac{\partial B}{\partial y}\right)dy\,dx
=\Big|_{x_0}^{x_1}\int_{y_0}^{y_1}A\,dy+\int_{x_0}^{x_1}\Big|_{y_0}^{y_1}\left(B-A\frac{dy}{dx}\right)dx,\tag{11}
\end{equation*}
wobei wir zur Abkürzung ausser von dem \textsc{Cauchy}schen Doppel-Substitutionszeichen
noch von der \textsc{Lindelöf}schen Bezeichnung
\begin{equation*}
\Big|_{y_0}^{y_1}A\frac{dy}{dx}=\frac{dy_1}{dx}\Big|^{y_1}A-\frac{dy_0}{dx}\Big|^{y_0}A\tag{12}
\end{equation*}
Gebrauch gemacht haben.

Die Anwendung der Formel (11) auf das in (9) vorkommende Doppelintegral
liefert für den einfachsten Fall, wo in dem Integral $J$ keine höheren
als die ersten Ableitungen der Funktion $z$ vorkommen, die folgende \textls[100]{zweite
Fundamentalformel von \textsc{Ohm}}\footnote{\textsc{Ohm}, loc. cit. 7), S. 306; bereits \textsc{Todhunter} hat auf die Formel (13) von \textsc{Ohm} aufmerksam gemacht, loc. cit. 5), S. 32.}:
\begin{equation*}
\begin{aligned}
\int_{x_0}^{x_1}\int_{y_0}^{y_1}\delta V\,dy\,dx
={}&\int_{x_0}^{x_1}\int_{y_0}^{y_1}\left(N-\frac{\partial P}{\partial x}-\frac{\partial Q}{\partial y}\right)\omega\,dy\,dx\\
&+\Big|_{x_0}^{x_1}\int_{y_0}^{y_1}P\omega\,dy+\int_{x_0}^{x_1}\Big|_{y_0}^{y_1}\left(Q-P\frac{dy}{dx}\right)\omega\,dx.
\end{aligned}\tag{13}
\end{equation*}
\par\smallskip\noindent{\small X 2 Abh. 5.}\hfill{\small 2}\par
\clearpage\SourceLeaf{467}
\ArithmeticHead{10}{L}{BOLZA, GAUSS UND DIE VARIATIONSRECHNUNG.}
\noindent\hspace*{\parindent}Der Ausdruck, der sich nunmehr durch Kombination von (9) und (13)
für $\delta J$ ergibt, hat die folgenden beiden Eigenschaften:

1) Die Integrationen und Substitutionen werden in allen Gliedern von $\delta J$
in derselben festen Reihenfolge ausgeführt.

2) In keinem Glied kommt unter einem Integralzeichen eine nach der
Integrationsvariabeln genommene partielle Ableitung von $\omega$ vor.

Das sind aber gerade die beiden charakteristischen Eigenschaften der
Normalform, auf die \textsc{Sarrus} in seiner Preisarbeit\footnote{\textit{Recherches sur le calcul des variations}, eingereicht 1842, jedoch erst 1848 veröffentlicht in den Mémoires présentés par divers savants à l’Académie des Sciences (Paris), Bd. X, S. 1—128.} die erste Variation eines
beliebigen $n$-fachen Integrals gebracht hat. \textsc{Ohm} hat also mit den beiden
Formeln (9) und (13) für den einfachsten Typus von Doppelintegralen die
\textsc{Sarrus}schen Resultate bereits vorweggenommen.

Damit war \textsc{Ohm} sehr nahe an die Ableitung der »Randbedingungen«
herangekommen; den letzten Schritt jedoch, die Aufstellung eines »Fundamentallemmas«,
das ihm gestattet hätte, alle Folgerungen aus dem Verschwinden
der ersten Variation zu ziehen, hat er nicht gemacht. Er hat zwar auch nach
dieser Richtung Untersuchungen angestellt; dieselben haben aber zu keinem
befriedigenden Ergebnis geführt\footnote{\textsc{Ohm}, loc. cit. 7), S. 75—80; der allgemeine Satz auf S. 79 ist unrichtig.}. — In seiner grossen Preisarbeit (1842) hat
dann \textsc{Sarrus} dieselbe Methode der Variation der Grenzen systematisch und
mit grösster Konsequenz auf beliebige vielfache Integrale ausgedehnt und die
ganze Theorie durch sein »allgemeines Fundamentallemma«\footnote{\textsc{Sarrus}, loc. cit. 10), S. 55—84.} erst zum vollen
Abschluss gebracht. Auch die gleichzeitige Preisarbeit von \textsc{Delaunay}\footnote{\textit{Mémoire sur le calcul des variations}, eingereicht 1842, veröffentlicht 1843 im Journal de l’École Polytechnique, Cahier 29, S. 37—120.} bedient
sich der Methode der Variation der Grenzen und ebenso die beiden
wichtigsten Lehrbücher der Variationsrechnung des neunzehnten Jahrhunderts,
\textsc{Jellett}\begingroup\renewcommand{\thefootnote}{13a)}\footnote{\textsc{J. H. Jellett}, \textit{An elementary treatise on the calculus of variations}, Dublin 1850; deutsch von \textsc{C. H. Schnuse}, Braunschweig 1860.}\endgroup\setcounter{footnote}{13}\ und \textsc{Lindelöf-Moigno}\textsuperscript{8a)}, die sich bezw. an \textsc{Delaunay} und \textsc{Sarrus}
anlehnen.
\par\clearpage\endgroup
